\documentclass[11pt]{article}

\usepackage[T1]{fontenc}
\usepackage[utf8]{inputenc}
\usepackage{lmodern}
\usepackage[margin=1in]{geometry}
\usepackage{microtype}
\usepackage{amsmath}
\usepackage{booktabs}
\usepackage{tabularx}
\usepackage{xcolor}
\usepackage[numbers,sort&compress]{natbib}
\usepackage[hidelinks]{hyperref}
\usepackage{cleveref}
\hypersetup{pdfauthor={dgxcoder}, pdftitle={ling-engine: An Exact, Bandwidth-Bound Inference Engine for One Model on One Unified-Memory Workstation}}

\newcommand{\code}[1]{\texttt{#1}}

\title{ling-engine: An Exact, Bandwidth-Bound Inference Engine\\for One Model on One Unified-Memory Workstation}

\author{dgxcoder\\ \texttt{dgxcoder@dreamference.ai}}

\date{October 2026}

\begin{document}
\maketitle

\begin{abstract}
A coding agent that runs on a single workstation spends almost all of its time waiting for its model.
On the NVIDIA DGX Spark (GB10, 128\,GB of unified LPDDR5X behind a 273\,GB/s bus), decoding a 27-billion-parameter
model is bound by memory bandwidth: every generated token streams the 17.6\,GB of weights once, which caps plain decoding
near 15 tokens/s however good the kernels are. We describe \emph{ling-engine}, a C++20/CUDA inference engine built for
exactly one model, Qwen3.8-27B in the mixed NVFP4/FP8 checkpoint that the Mightling agent serves, and for exactly one
workload, that agent's recorded sessions. The engine keeps two invariants that general-purpose servers do not: greedy
speculative output is token for token identical to plain decoding, and the recurrent state after a speculative run is bit
for bit the state of a plain run, which lets every optimization be accepted or rejected by an equality test rather than a
quality benchmark. Four measured milestones took the engine from 33.7 to 49.6 tokens/s of single-stream decode on
replayed agent sessions, against 47--48 tokens/s for the production server (SGLang with the same drafter) on the same
machine; prefill runs 1.1--1.3$\times$ production's; two agents served at once decode at 34.3 tokens/s per stream against
production's 33.9. We report what did not pay, among them a context-lookup draft branch and a third pipeline stage, and
one methodological finding: unseeded replays of sampled sessions carry $\pm5\%$ of acceptance noise, enough to hide or
invent every gain reported here, so every comparison is now made on seeded replays with identical output.
\end{abstract}

\section{Introduction}

Mightling is a local, air-gapped coding agent for a single DGX Spark. Its agent, \code{ling}, is a branded build of a
vendor command-line agent, and its model server is SGLang running Qwen3.8-27B with the DFlash2 block
drafter~\cite{chen2026dflash,dflash2drafter}. Measured over 168 recorded agent sessions (11,431 model requests), 92\% of a
session's wall time is spent in the model; a request writes a median of 98 and a mean of 177 tokens, 71\% of them
tool-call arguments; 46\% of what the agent writes is copied from its own context; and every new session re-reads a
system prompt and tool schema of about 11--15K tokens that all sessions share. Engine speed therefore turns almost one for
one into agent speed, and the relevant benchmark is a replay of those sessions, not a chat benchmark.

A general server spends part of each decode step on work that is not weight traffic. ling-engine is the opposite bet: a
compiler and a runtime for one model and one checkpoint, where every kernel is specialized to a matrix shape, every decode
step is a captured CUDA graph, and the output is provably the same as plain decoding. This paper reports its design, the
discipline that keeps it exact, four milestones of measurements on the same machine as production, and what did not work.

\paragraph{Contributions.}
(1) A single-model engine whose speculative decoding is exact under two checkable rules, row invariance and commit by
replay (\cref{sec:exact}).
(2) Measurements on one machine of where a speculative step's time goes, before and after each change
(\cref{sec:results}).
(3) Negative results on the context-lookup draft branch, pipeline depth, draft-block width under batching and FP8 KV on
this hardware (\cref{sec:negative}).
(4) The observation that replay benchmarks of sampled sessions must be seeded to be comparable at all
(\cref{sec:noise}).

\section{Hardware and model}

\paragraph{The bus is the budget.}
A read-only streaming kernel sustains 262--264\,GB/s on the GB10, 96--97\% of the 273\,GB/s peak; four SMs are enough
to reach 98\% of that, and DRAM latency is about 400\,ns. CPU memory traffic shares the same pins byte for byte: two CPU
threads streaming their own buffers take 27\,GB/s from the GPU's read stream. The engine's targets are stated against
262\,GB/s, and anything co-resident with a model server, such as an indexer, is budgeted as bus traffic.

\paragraph{The model.}
Qwen3.8-27B is a dense hybrid of 64 layers in 16 repeats of three Gated DeltaNet blocks followed by one gated
full-attention block, each with its own SwiGLU FFN (hidden 5,120, intermediate 17,408), a 248,320-token vocabulary and a
262,144-token native context. The checkpoint Mightling serves is mixed precision: the FFNs and the LM head are NVFP4
(one FP8 scale per 16 values), every DeltaNet and attention projection is FP8. One pass over the target weights reads
17.6\,GB (\cref{tab:bytes}); moving the FP8 projections to NVFP4 would save 18\%, which the checkpoint's makers chose
not to do and which the engine therefore does not do either. Per-sequence state is small by comparison: a BF16 KV cache
of 64\,KB per token over the 16 attention layers and a DeltaNet state of 151\,MB in FP32 that does not grow with context.

\begin{table}[t]
\centering\small
\begin{tabular}{lrr}
\toprule
Tensor group & Precision & GB per pass \\
\midrule
FFN gate, up, down (64 layers) & NVFP4 & 9.63 \\
DeltaNet in- and out-projections (48 layers) & FP8 & 5.54 \\
Attention q, k, v, o (16 layers) & FP8 & 1.68 \\
LM head & NVFP4 & 0.72 \\
$\beta/\alpha$ projections, norms, conv & BF16 & 0.05 \\
\midrule
Target weights per pass & & 17.6 \\
Drafter layers & NVFP4 & 1.55 \\
\bottomrule
\end{tabular}
\caption{What one decode or verify pass streams. At 262\,GB/s the target weights alone take 67\,ms.}
\label{tab:bytes}
\end{table}

\paragraph{The drafter.}
DFlash2 drafts a block of $B$ tokens in one pass from the target's residual stream after five of its layers, and a
candidate selector walks a $16\times16$ lattice per position to one chain, which the target verifies in a single pass of
$B$ rows. Production serves $B=16$; the engine runs the drafter exactly as production's code does and reproduces its
acceptance, 4.98 against 4.95 tokens per step on the replay.

\section{Design}

\paragraph{Two programs.}
\code{ling-compile} is an offline step that fixes the precision map, folds the norms, lays the weights out on disk in the
order the kernels read them, and instantiates one fused kernel per matrix shape. \code{ling-serve} is a resident runtime
that keeps every session's KV cache and DeltaNet state on the GPU and speaks the agent's API, so the agent cannot tell
which server answers.

\paragraph{The rows path.}
Every pass of 32 rows or fewer (a decode step, a verify, a short prompt) goes through one family of kernels:
\code{stream\_gemm}, a tensor-core weight-streaming GEMM for 1--32 rows that converts NVFP4 and FP8 weights to FP16 in
registers exactly, with FP32 accumulation and split-K chosen from the matrix shape; a narrow BF16 kernel for the DeltaNet
$\beta/\alpha$ projections; and tensor-core attention over fixed 512-key ranges. Weights are retiled at load into the order
the kernel reads them; the retiling alone gave 5--6\% end to end.

\paragraph{Prefill.}
Prompts of more than 32 rows go through a separate GEMM on block-scaled tensor cores (\code{mma.sync kind::mxf4nvf4} for
NVFP4, \code{kind::f8f6f4} for FP8) over the same tiled weights, with activations quantized exactly as production's
server quantizes them. The engine's own kernel was kept over CUTLASS~4.5.1 because it is faster on one of the two FFN
shapes and slower on the other, and one tiled copy of the weights serves both paths.

\paragraph{One graph per step.}
A speculative step has two data-dependent decisions, the lattice walk after the draft and the accept rule after the
verify, and both already ended in a synchronization. The three GPU phases between them, about 200, 1,300 and 100 launches,
are each captured once into a CUDA graph and replayed every step; positions and tokens travel through device memory and
pinned host buffers at fixed addresses, and the commit has one graph per accepted count. Programmatic dependent launch lets
each kernel on the rows path start as its predecessor drains, under one rule: every kernel waits on the grid dependency
before its first read of a predecessor's output and before its first write.

\paragraph{Sessions and batches of two.}
The runtime holds a fixed number of resident sessions, each with its KV caches, DeltaNet and conv state, an end-of-prompt
snapshot and a 2,048-position ring for the drafter's context; a request binds to the session whose history its prompt
extends, and prefix checkpoints at shared-prefix boundaries let a new session resume after the system message every
session shares. Two sessions decode in one pass: their verify rows are concatenated ($V_1+V_2\le32$), the per-row
kernels take the second sequence unchanged, and the stateful kernels (attention, DeltaNet recurrence, the drafter's cache)
run once over a device-side descriptor of both sequences. A fresh long prompt arriving while another session decodes is
interleaved in 256-token chunks, so the decoding stream's step period stays within about twice its normal step.

\section{Exactness}
\label{sec:exact}

Speculative decoding preserves the output distribution in theory~\cite{leviathan2023speculative}; the engine makes it hold
in the bits, so that every change can be judged by equality. Two design rules carry it.

\paragraph{Row invariance.}
Every kernel on the rows path computes a row bit for bit the same whatever the other rows in the pass are: the GEMM uses a
fixed K-chunk order per warp, a fixed-order warp reduction and a split-K never chosen from the row count; attention runs
over fixed 512-key ranges at fixed positions, where a tile past a row's last key leaves the online-softmax state exactly
unchanged. A verify of 16 rows therefore computes for row $i$ exactly what plain decoding would compute at that position.

\paragraph{Commit by replay.}
A verify reads the DeltaNet state and does not write it; it keeps each layer's projected inputs. The commit replays the
accepted rows through the same convolution and recurrence kernels, whose per-token arithmetic does not depend on how many
tokens one call covers. The state after a speculative run is thus the state of a plain run over the same tokens.

\paragraph{The gates.}
Every build runs \code{--spec-check} (each prompt decoded plainly and speculatively, tokens compared, then the speculative
tokens replayed through plain decoding and every DeltaNet and conv state compared bit for bit), \code{--prefix-check} (a
prompt resumed from a checkpoint reaches the state a prefill from scratch does), \code{--graph-check} (each prompt with and
without graphs, greedy and seeded, requiring identical tokens, state and per-step acceptance), and from M4 on
\code{--interleave-check} and \code{--batch-check} (each sequence of a pair equal to itself alone, in both row orders,
with either one prefilling beside the other). Sampling is exact rejection sampling against the selector's 16-way draft
distribution, tested by chi-square on a toy model. An optimization merges only when every gate passes, the output is the
same as the previous main's, and it is no slower at one stream and at two. The graph check found one real defect: the first
draft after every prompt had been running the drafter's projections through the prefill path.

\section{Results}
\label{sec:results}

All measurements are on one development GB10 with production's SGLang image and flags as the reference, on a replay of
recorded agent sessions (M0's main set: 12 sessions $\times$ 12 requests, production sampling, single stream unless
stated). Recorded sessions never leave the machine; only aggregates are reported.

\paragraph{Where production's step goes.}
Production decodes 47 tokens/s at a 105\,ms step and 4.95 accepted tokens per step. An Nsight trace splits the step into
86\,ms of target GEMMs at 204\,GB/s (78\% of the bus), 8.6\,ms of drafter, 7.5\,ms of DeltaNet verify, 6.1\,ms of
attention over a BF16 KV cache and 3\,ms of commit, KV append and sampling; the GPU idles only about 1\,ms. The spec had
assumed ``30 tokens/s at 160\,ms'', which turned out to be production under two concurrent sessions on a loaded host. A
side result of the same measurement: production itself gains 7\% on one stream from 12 draft tokens instead of 16
(48.3 against 45.1 tokens/s), exactly, which Mightling adopted.

\paragraph{Decode across milestones.}
\Cref{tab:decode} follows the single-stream step. M1 reached 97\% of production's decode speed with the verify's weight
GEMMs already faster than production's (81.7 against 86.3\,ms for the 17.6\,GB) but a longer step elsewhere. M3 took
4.7\,ms off with graphs ($-3.0$), PDL ($-0.7$) and an L2 prefetch in the verify attention ($-1.0$ to $-1.5$). M4's
session and batch work left the one-row step unchanged within noise. The step-savings block then measured each kernel
change on its own: B16, two 16-row tiles per block on the narrow matrices with two chunks in flight per warp, and the
commit's DeltaNet replay moved to a second stream beside the next draft.

\begin{table}[t]
\centering\footnotesize
\begin{tabular}{llrrr}
\toprule
Build & What changed & ms/step & tok/step & tok/s \\
\midrule
Production (SGLang, 16 drafts) & reference & 105.5 & 4.95 & 47.1 \\
First M1 build & speculation working & 146 & & 33.7 \\
M1 & retiled weights, streaming GEMMs & 114.5 & 5.21 & 45.5 \\
M3 & graphs, PDL, attention prefetch & 110.5 & & 43.5 \\
M4 stage 4 & sessions and batches; one-row path untouched & 104.2 & 5.17 & 49.6 \\
B16 (seeded) & weight stream & 102.6 & & 48.6 \\
Commit overlap (seeded) & DeltaNet replay on a second stream & 101.4 & 4.99 & 49.2 \\
\bottomrule
\end{tabular}
\caption{Single-stream decode on the replayed agent sessions, server-measured. Rows before ``seeded'' were replayed
without a fixed sampling seed and differ in acceptance between runs (\cref{sec:noise}); the two seeded rows produced
the same 17,527 tokens in the same 3,514 passes.}
\label{tab:decode}
\end{table}

\paragraph{Prefill and request latency.}
v0's prefill dequantized every weight per 2,048-token chunk and ran cuBLAS, so a warm request paid 1.7\,s before its
first token. The block-scaled GEMM runs prefill at 2,637 tokens/s on a 1,061-token prompt, 2,584 on 8,313 and 2,106 on
34,778, against production's 2,142--2,462, 1,934 and 1,683, and 3.6--4.8$\times$ M1's. On the replay a request then
takes 4.30--4.43\,s on average against M1's 8.14 and production's 4.18--4.25; a warm request's first token comes after
0.40\,s median against production's 0.49--0.50; a session's first request after 4.4--5.2\,s against 8.5--10.1, because a
checkpoint after the shared system message replaces its prefill. Prompts are token-identical to production's on all
396 replayed requests tested, after three rendering differences were found and copied, one of them production's own
tokenizer behaviour on combining marks.

\paragraph{Two streams.}
The M4 gate was ``two agent tasks served at once at no less than production's two-stream numbers'': 33.9 tokens/s per
stream and a 145\,ms step period. Stage 2's interleaved scheduler alone gave 22.1 tokens/s per stream (the two streams
took turns); stage 3's batched verify 30.7--31.2 at a 157\,ms period with 81--83\% of passes batched; stage 4, which runs
the drafter over both sequences' rows in one pass and every state kernel in one launch, 33.4 pooled at 151.9\,ms. On
seeded replays B16 took the per-stream period from 154.2 to 147.5\,ms (32.6 to 34.1 tokens/s) and the commit overlap
from 147.8 to 146.3 (34.0 to 34.3), both with identical output, so the engine now meets the two-stream gate. The batched
step itself is about 122\,ms for two sequences, against 102 for one: the second sequence costs its attention and state
traffic, not a second weight read.

\section{What did not pay}
\label{sec:negative}

\paragraph{The context-lookup draft branch.}
The spec's second lever was a lookup of the continuation of the last three tokens' most recent earlier occurrence, verified
in place of the drafter's chain on a long match, on the argument that 46\% of agent output is copied. Scored in shadow on
every step, it lost to the drafter in every match-length bucket, even on verbatim copies of 32 tokens or more (6.9
against 9.9 accepted), and the best of both per step was only $+4\%$; the simulated $\times1.2$ had assumed independent
successes. A later statistical study over 972 sessions (16.25M output tokens) gave the mechanism: long copies exist
(45.9\% of tokens in runs of 16 or more, 7.8\% in runs of 128 or more) but the drafter already accepts 8--11 tokens per
step on exactly those spans, so a realistic copy proposer adds 6.8\% in tokens per step at 64 verify rows, and a copy
step costs about as much as a drafter step. The idea survives only as a gated proposer that takes the copy when it beats
the drafter, under a stricter whole-session bar.

\paragraph{A wider or deeper pipeline.}
Block 12 shortens the engine's step by 2.8\,ms but loses 10\% of acceptance (45.8 against 49.8 tokens/s), the opposite of
production's experience; the engine keeps 16. A third \code{cp.async} stage in the verify attention (B07) was exact and
slower at two streams (122.3 to 124.1\,ms per batched step) because the kernel already read its KV at 219\,GB/s against a
plain read's 225. A per-sequence draft block of 16 under batching (B31) was worth about 1\% and was closed. FP8 KV gave no
step gain on this hardware, and a wide variant of the streaming GEMM for prefill ran at a third of the speed because every
16-row weight tile re-read all 128 activation rows from L2.

\paragraph{What was free.}
Two sessions in one pass cost 17\,ms over one; the drafter over both sequences' rows saved 7\,ms (its weights read once);
one launch per state kernel and graphs for the batch path saved nothing measurable, since at 24 rows the verify is its
weight stream plus the second sequence's attention.

\section{Measure the noise first}
\label{sec:noise}

Until October 14 the replay harness sent no sampling seed, so each request sampled from a fresh random seed and two runs
of ``the same'' replay were two workloads. The run-to-run spread in accepted tokens per step was $\pm5\%$ (4.87 against
5.28 on one pair), which moves tokens/s by the same amount while moving the step period much less. That spread is the
size of every gain in \cref{tab:decode}, so B16's first two-stream pair read as a loss in tokens/s (33.9 to 33.1) while
its step period had shortened by 8.8\,ms with every output bit unchanged. The harness now fixes a seed per request from
the run, the session and the request index; a comparison is reported only when the two runs produced the same output
tokens in the same number of passes, which makes tokens/s a measurement of the step rather than of the dice. Even seeded,
about one request in 144 at two streams can still differ, because resuming a session from a slot rather than re-prefilling
changes the drafter's context and hence its proposals; greedy output never changes. Earlier comparisons, including
production's own 47--48 tokens/s rows, carry the unseeded noise and are reported as such.

\section{Status and plan}

The engine serves the agent's full API (chat, the Responses conversion, reasoning split, streaming, typed tool calls,
completions from ids) and the production server's metric names, so the same replay harness and the same agent run
against either. Greedy top-1 agreement with production is 98.6\% and top-5 agreement 100\%; the mean KL against
production's own run-to-run noise is within that noise. The engine has not yet replaced production: the owner's bar is a
benchmark night, two agent tasks at once on real SWE-bench instances, at no less than production's numbers, and the
remaining single-step savings (a context-aware sparse drafter head, 6-bit scale codes, and re-quantizing the drafter
until it accepts as its BF16 original does) are measured one at a time against the main of their day. After that:
structured output and images, which the agent needs; the spec's tree verify; and a repository-aware copy proposer under
its gate. Tensor parallelism across two Sparks is last, after every one-node improvement is settled.

\section{Limitations}

Every number is from one machine, one model and one checkpoint, on a replay of one agent's recorded sessions; the engine
is not a general server and does not try to be. Rows measured before the seed fix carry acceptance noise of the size of the
reported differences, and the production rows were taken the same way. The decode measurements are server-side and warm;
the agent's end-to-end gain depends on the request mix, which is dominated by short decodes. Nothing here is a leaderboard
score.

\paragraph{Reproducibility.}
The engine, its specification, the milestone reports and the bench scripts are in the ling-engine repository; the
recorded sessions are private and the harness skips the tests that need them.

\bibliographystyle{plainnat}
\bibliography{references}

\end{document}
